\chapter{GPS- und Kinematikberechnung}
\label{ch:gps-kinematik}

Die GPS-Eingabe ist eine semikolonseparierte CSV mit den Spalten \texttt{time}, \texttt{lat}, \texttt{lon}, \texttt{ele} und \texttt{temperature}. Die Zeitstempel liest \texttt{pandas.to\_datetime} als UTC ein \cite{pandas_docs}. Die Zeilen müssen bereits zeitlich sortiert sein, die Berechnung nutzt dann nur aufeinanderfolgende Punkte.

\section{Zeit und Distanz}

Für zwei aufeinanderfolgende Punkte $i-1$ und $i$ gilt im Code:
\begin{equation}
\Delta t_i = t_i - t_{i-1}
\end{equation}
\begin{equation}
\Delta s_i = 2 R \arcsin\left(\sqrt{a}\right)
\end{equation}
mit
\begin{equation}
 a = \sin^2\!\left(\frac{\Delta \varphi}{2}\right) + \cos(\varphi_{i-1})\cos(\varphi_i)\sin^2\!\left(\frac{\Delta \lambda}{2}\right)
\end{equation}
und $R = \SI{6371000}{\metre}$ als mittlerem Erdradius. Diese Haversine-Formel ist für die Distanzberechnung vorgegeben und wird direkt in \texttt{GPSTrack\_distanz\_haversine()} umgesetzt.

Aus Distanz und Zeit folgt die Rohgeschwindigkeit:
\begin{equation}
 v_i = \frac{\Delta s_i}{\Delta t_i}
\end{equation}

Ist \(\Delta t_i \le 0\), setzt das Projekt die Rohgeschwindigkeit auf \SI{0}{\metre\per\second} und markiert das Intervall als ungültig. So vermeidet die Berechnung eine Division durch null, ohne Datenpunkte zu verlieren.

\section{Steigung, Kurs und Kompassrichtung}

Die Höhendifferenz wird als
\begin{equation}
\Delta h_i = h_i - h_{i-1}
\end{equation}
verwendet. Daraus entsteht die Steigung in Grad als
\begin{equation}
\varphi_i = \arctan\!\left(\frac{\Delta h_i}{\Delta s_i}\right)
\end{equation}
falls $\Delta s_i > 0$ ist, sonst $0$. 
Der Kurswinkel (englisch \emph{initial bearing}) folgt aus Breiten- und Längengrad
zweier aufeinanderfolgender Punkte über die Formel
\begin{equation}
    \theta_i = \operatorname{atan2}\Big(
        \sin(\Delta\lambda)\cos(\varphi_i),\;
        \cos(\varphi_{i-1})\sin(\varphi_i) - \sin(\varphi_{i-1})\cos(\varphi_i)\cos(\Delta\lambda)
    \Big)
    \label{eq:kurswinkel}
\end{equation}
mit $\Delta\lambda = \lambda_i - \lambda_{i-1}$. Da $\operatorname{atan2}$ Werte im Bereich
$(-180\si{\degree}, 180\si{\degree}]$ liefert, wird das Ergebnis anschlie{\ss}end über
$\theta_i \bmod 360\si{\degree}$ auf $[0\si{\degree}, 360\si{\degree})$ normalisiert.

Aus dem normalisierten Kurswinkel ergibt sich die Kompassrichtung durch Rundung auf
den nächstgelegenen $45\si{\degree}$-Sektor:
\begin{equation}
    \text{Sektorindex} = \operatorname{round}\!\left(\frac{\theta_i}{45\si{\degree}}\right) \bmod 8
    \label{eq:kompasssektor}
\end{equation}
Die acht Sektoren $0$ bis $7$ entsprechen dabei den Richtungen N, NO, O, SO, S, SW, W
und NW, beginnend bei $0\si{\degree}$ (Norden) und im Uhrzeigersinn fortlaufend. Ein
Kurswinkel von z.\,B. $80\si{\degree}$ ergibt so Sektor $2$, also O.

\section{Beschleunigung}

Die Beschleunigung entsteht nicht aus einer einfachen Differenz zweier aufeinanderfolgender Geschwindigkeitswerte, sondern segmentweise aus der aktiven Geschwindigkeitskurve. Der Grund: Bei ungleichmäßigen Zeitabständen -- wie sie bei GPS-Aufzeichnungen üblich sind -- würde eine reine Vorwärtsdifferenz das Ergebnis unnötig empfindlich gegenüber genau dem Punkt machen, an dem sie ausgewertet wird. \texttt{beschleunigung\_aus\_geschwindigkeit()} verwendet deshalb eine zentrale Differenz: Für einen Punkt $i$ wird die Geschwindigkeitsänderung zwischen dem Vorgänger- und dem Nachfolgerpunkt durch die Zeit zwischen den beiden Intervallmitten geteilt. Das mittelt den Einfluss beider Nachbarn und reagiert dadurch weniger empfindlich auf einen einzelnen verrauschten Messwert als eine reine Vorwärts- oder Rückwärtsdifferenz.

An den Segmentgrenzen gibt es keinen Vorgänger bzw. Nachfolger im selben Segment; dort setzt die Funktion die Beschleunigung am ersten Punkt auf null und verwendet am letzten Punkt eine einseitige (rückwärtige) Differenz. Segmente werden dabei nicht überschritten: Eine Beschleunigung wird nie über eine grosse Zeitlücke hinweg berechnet, weil das einen physikalisch bedeutungslosen Wert ergeben würde (z.\,B. eine sehr kleine Beschleunigung über eine mehrminütige Aufzeichnungspause).

\section{Wichtige DataFrame-Spalten}

Diese Spalten bilden die Grundlage für die weiteren Kapitel:
\texttt{dt\_s}, \texttt{distanz\_m}, \texttt{steigung\_grad}, \texttt{kurs\_grad}, \texttt{kompassrichtung}, \texttt{geschwindigkeit\_roh\_ms}, \texttt{geschwindigkeit\_geglaettet\_ms}, \texttt{geschwindigkeit\_ms}, \texttt{beschleunigung\_roh\_ms2}, \texttt{beschleunigung\_geglaettet\_ms2}, \texttt{beschleunigung\_ms2} sowie die Filterspalten der Glättung. Der Simulator nutzt später nur \texttt{geschwindigkeit\_ms} und \texttt{beschleunigung\_ms2}.
